% SOURCES. Binding final4/GOAL_alt.md; independent symbol differentiation.
% [X] Named agreeing routes: this author's Heisenberg-velocity/band projection
% and the icloud lane's Bloch-group derivative in E7--E8,E18 of
% <internal>/strand_E_declared_equations.tex,
% read in full including SOURCES after this independent calculation was saved.
% The shared exact N Q_i/E+v_i cos(2ap_i) formula supplies the displayed Taylor
% and packet-average corollaries. No agreement on an unspecified collapse
% history or finite-spacing geometric evolution is claimed; counterpart
% overall review remains pending. Predecessor lane_z/B11-007's original
% g30b source was read: its cos(2ap_i) dressing is precisely the flow velocity
% derivative here. Its high-band discrepancy is retained, not removed.
% Derivation: for the positive band, partial_i E=Q_i/E exactly.
% Heisenberg velocity of the full spinor is N alpha_i cos(ap_i)+v_i cos(2ap_i);
% its positive-band expectation gives the displayed formula.
% A finite packet is not a momentum eigenstate. If its momentum distribution
% is inversion symmetric around a valley, <Q_i/E>=0, but
% <cos(2ak_i)>=1-2a^2<k_i^2>+O(a^4<k_i^4>).
% In a slowly varying background the formula is its local geometrical-optics
% velocity, subject to packet-size/background-scale and interband errors.
% For m=0 at a valley E=0 the positive-band velocity has no direction-independent
% limit: no massive rest-packet theorem should be claimed for that point.

\paragraph{Comoving matter [X].}
For a positive-energy packet in a uniform configuration, the exact band
velocity is
\[
 \dot x_i=N{Q_i(p)\over E(p)}+v_i\cos(2ap_i).
\]
For $m>0$, the zero relative momentum at every valley gives $\dot x=v$.
At finite $k=p-\nu\pi/a$, put $E_0=\sqrt{m^2+k^2}$; then
\[
 \dot x_i=v_i+{Nk_i\over E_0}
 -a^2\left(2v_i k_i^2+{2Nk_i^3\over3E_0}
       -{Nk_i\sum_jk_j^4\over6E_0^3}\right)+O(a^4).
\]
A finite packet centered symmetrically at the valley instead has mean
velocity $v_i[1-2a^2\langle k_i^2\rangle+O(a^4\langle k_i^4\rangle)]$.
In a slowly varying $N=1$ flow, the ray equation
$\dot k=-\nabla_x\omega$ preserves $k=0$, so massive dust follows $v$.
The massless band has no rest point; its directed rays move at $v+N\boldsymbol n$
in the continuum limit.
